\chapter{Introduction}
\label{ch:introduction}

\section{Domain Context}

Organizations and career platforms increasingly rely on natural language processing to interpret resumes, job postings, and professional profiles.
Skill mentions in these documents are rarely presented as normalized database entries; they appear as spans within sentences, often with ambiguous boundaries and mixed hard and soft skill types~\cite{zhang2022skillspan}.
Accurate extraction of such spans is a prerequisite for skill-based job matching, labor-market analytics, and personalized learning recommendations.

\section{Problem Definition}

This thesis addresses \emph{skill span extraction} from English employment-related text: given an input text segment, the system must identify contiguous substrings that denote explicit hard or soft skills and assign each span a label from a fixed schema.
The task follows the SkillSpan formulation, which distinguishes \textsc{Hard\_Skill} and \textsc{Soft\_Skill} annotations on sentence-level examples drawn from job postings~\cite{zhang2022skillspan}.

Beyond extraction, the Stage-2 SkillLLM\_v2 system targets a broader \emph{skill intelligence loop}: user profile construction, taxonomy normalization, job matching, skill-gap detection, and course recommendation.
Only extraction and wrapper components were evaluated on SkillSpan gold labels; downstream matching metrics lacked labeled evaluation data in the repository artifacts inspected for this thesis.

\section{Motivation}

Prior work shows that repurposing general-purpose large language models (LLMs) through supervised adaptation can improve skill extraction relative to lighter baselines~\cite{herandi2024skillllm}.
Open-weight models introduce practical trade-offs among model capacity, inference cost, and hardware feasibility.
The repository records both a multi-model benchmark for model selection and a subsequent system integration effort that treats a single selected extractor as the non-regression anchor for wrapper components.

\section{Research Questions}

The experiments completed in the repository support the following research questions:

\begin{description}
  \item[RQ1] How do contemporary open-weight instruction-tuned LLMs compare when fine-tuned under a shared SkillSpan protocol with identical splits and metrics?
  \item[RQ2] Among models feasible on 2$\times$ NVIDIA RTX 3090 (24\,GB) GPUs, which backbone yields the strongest span extraction quality?
  \item[RQ3] How do lightweight encoder baselines (BERT token classification, zero-shot GLiNER) compare with a fine-tuned generative Mistral extractor on SkillSpan?
  \item[RQ4] Can a guarded hybrid wrapper---combining Mistral extraction with auxiliary candidates, boundary refinement, and taxonomy linking---improve or preserve span metrics while enabling downstream normalization?
  \item[RQ5] What systematic error modes remain after fine-tuning and wrapper processing?
\end{description}

RQ4's downstream matching and recommendation claims are limited to implemented pipeline behavior and linking coverage; nDCG and precision@k were not computed because labeled job/course evaluation sets were absent from verified artifacts.

\section{Contributions}

This thesis documents the following evidence-backed contributions:

\begin{enumerate}
  \item A reproducible multi-model SkillSpan benchmark pipeline (\texttt{SkillLLM\_v2}) comparing five open-weight LLMs with shared preprocessing, LoRA/QLoRA training, JSON generative decoding, span alignment, and multi-metric evaluation.
  \item Empirical model selection showing Mistral Small 3.1 24B as the strongest completed benchmark model among feasible configurations, with documented failure of Gemma 4 26B on available hardware.
  \item SkillLLM\_v2 Stage-2: a context-aware graph-augmented architecture with hybrid candidate extraction, non-regression gating, taxonomy linking, and NetworkX-based relational reasoning for job and course recommendation.
  \item Verified comparative metrics on SkillSpan test data for BERT, GLiNER, raw Mistral LoRA, and wrapper variants, including wrapper regression analysis against the Mistral anchor.
\end{enumerate}

\section{Thesis Organization}

Chapter~\ref{ch:related} reviews related work.
Chapter~\ref{ch:methodology} formalizes the task and evaluation methodology.
Chapter~\ref{ch:dataset} documents datasets and preprocessing.
Chapter~\ref{ch:architecture} describes the SkillLLM\_v2 architecture.
Chapter~\ref{ch:baselines} presents baseline systems.
Chapter~\ref{ch:experiments} details experimental setup.
Chapter~\ref{ch:results} reports results.
Chapters~\ref{ch:discussion}--\ref{ch:error} analyze findings and errors.
Chapter~\ref{ch:reproducibility} summarizes reproducibility.
Chapters~\ref{ch:limitations}--\ref{ch:conclusion} discuss limitations, future work, and conclusions.

\section{Formatting Note}

This document was prepared for CSC 898 (Master's Thesis) at San Francisco State University.
Graduate Studies requires department-approved formatting, 1-inch margins, accessible headings, and figure/table captions~\cite{sfsugrad2026thesis,sfcs2026thesis}.
No department-provided \LaTeX{} template was present in the repository; final format approval by the Graduate Program Specialist remains required.
